\SourcePage{167}
\SGAchapter{XIV}{深度与平展上同调中的 Lefschetz 定理}{Profondeur et théorèmes de Lefschetz en cohomologie étale}
\begin{center}M. Raynaud 夫人\footnote{根据 A. Grothendieck 未发表的笔记。}\end{center}
第1节定义一种“\Term{etale-depth}{平展深度}”，它是在\Term{etale-cohomology}{平展上同调}中与第 III 讲所研究的\Term{coherent-sheaf}{凝聚层}\Term{cohomology}{上同调}中的\Term{depth}{深度}相对应的概念。经过一部分技术性准备后，我们在第4节证明若干“Lefschetz 定理”，其中的\Term{center}{中心}定理是4.2。设 $X$ 为\Term{scheme}{概形}，$Y$ 为 $X$ 的\Term{closed-subset}{闭子集}，$U$ 为补\Term{open-subset}{开集} $X-Y$，$F$ 为 $X$ 的\Term{etale-topology}{平展拓扑}上的\Term{abelian-sheaf}{阿贝尔层}。一般说来，Lefschetz 定理的目的是说明：若 $F$ 在 $U$ 上满足某些可以用 $U$ 各点的\Term{etale-depth}{平展深度}表述的\Term{local-scheme}{局部}条件，则在关于 $U$ 的某些附加整体条件下（例如 $U$ 为\Term{affine}{仿射}），自然映射
\[H^i(X,F)\longrightarrow H^i(Y,F|_Y)\]
对 $i<n$ 为\Term{isomorphism}{同构}，其中 $n$ 是一个明确给出的整数。取 $F$ 为\Term{constant-sheaf}{常值层}，便得到 $\pi_0(X)$ 等于 $\pi_0(Y)$ 的条件，以及 $X$ 与 $Y$ 的\Term{fundamental-group}{基本群}的\Term{abelianization}{阿贝尔化}相同的条件。第5节引入“\Term{geometric-depth}{几何深度}”，由此给出 Lefschetz 定理的一些有用特例（5.7）。最后，第6节指出若干猜想，特别涉及所得定理的“非交换”变体。
\SGAsection{1}{上同调深度与同伦深度}
\textbf{1.0。}固定如下记号。设 $X$ 为\Term{scheme}{概形}\footnote{依照新的术语（参看 EGA I 再版），此处把原先称为“\Term{prescheme}{预概形}”的\Term{object}{对象}称为“\Term{scheme}{概形}”，把原先称为“\Term{scheme}{概形}”的\Term{object}{对象}称为“\Term{separated-scheme}{分离概形}”。}，$Y$ 为其\Term{closed-subset}{闭子集}，$U$ 为补\Term{open-subset}{开集}，$i:Y=X-U\to X$ 为\Term{canonical}{典范}\Term{immersion}{浸入}。记 $\underline\Gamma_Y$ 为把 $X$ 上的\Term{abelian-sheaf}{阿贝尔层}映为“\Term{support}{支集}含于 $Y$ 的\Term{section}{截面}层”的\Term{functor}{函子}，即 $\underline\Gamma_Y=i_*i^!$（参看 SGA 4 IV 3.8 与 VIII 6.6），
\SourcePage{168}
记 $\Gamma_Y$ 为\Term{functor}{函子} $\Gamma\cdot\underline\Gamma_Y$（其中 $\Gamma$ 是“\Term{global-section}{整体截面}”\Term{functor}{函子}）。考虑\Term{derived-category}{导出范畴} $D^+(X)$ 以及 $\underline\Gamma_Y$（分别 $\Gamma_Y$）的\Term{derived-functor}{导出函子} $R\underline\Gamma_Y$（分别 $R\Gamma_Y$）（参看[3]）。给定 $X$ 上次数有下界的\Term{abelian-sheaf}{阿贝尔层}\Term{complex}{复形} $F$，可以把它视为 $D^+(X)$ 的\Term{object}{对象}；以 $\mathcal H_Y^p(F)$ 记 $R\underline\Gamma_Y(F)$ 的第 $p$ 个\Term{cohomology-sheaf}{上同调层}，以 $H_Y^p(X,F)$ 记 $R\Gamma_Y(F)$ 的第 $p$ 个\Term{cohomology-group}{上同调群}。（SGA 4 V 4.3 和4.4）的结果显然推广到 $\mathcal H_Y^p(F)$ 与 $H_Y^p(X,F)$。
\begin{SGAresult}{命题}{1.1}
设 $X$ 为\Term{scheme}{概形}，$Y$ 为\Term{closed-subset}{闭子集}，$U$ 为补\Term{open-subset}{开集}，$i:U\to X$ 为\Term{canonical}{典范}\Term{immersion}{浸入}。$F$ 表示 $X$ 上的\Term{sheaf-of-sets}{集合层}、\Term{sheaf-of-groups}{群层}，或次数有下界的\Term{abelian-sheaf}{阿贝尔层}\Term{complex}{复形}。固定如下记号：若 $X'\to X$ 是\Term{morphism}{态射}，以 $U'$ 与 $F'$ 表示 $U$ 与 $F$ 在 $X'$ 上的\Term{inverse-image}{逆像}；若 $\bar y$ 为 $X$ 的\Term{geometric-point}{几何点}，以 $\overline X$ 表示 $X$ 在 $\bar y$ 处的\Term{strict-localization}{严格局部化}，以 $\overline U$ 与 $\overline F$ 表示 $U$ 与 $F$ 在 $\overline X$ 上的\Term{inverse-image}{逆像}。

$1^\circ)$ 设 $F$ 为 $X$ 上的\Term{sheaf-of-sets}{集合层}，$n\leq2$ 为整数，则下列条件\Term{equivalence}{等价}：

(i) \Term{canonical}{典范}\Term{morphism}{态射} $F\to i_*i^*F$ 在 $n\geq1$ 时为\Term{injectivity}{单射}，在 $n\geq2$ 时为\Term{bijection}{双射}。

(ii) 对每个在 $X$ 上\Term{etale}{平展}的\Term{scheme}{概形} $X'$，\Term{canonical}{典范}\Term{morphism}{态射}
\[H^0(X',F')\longrightarrow H^0(U',F')\]
在 $n\geq1$ 时为\Term{injectivity}{单射}，在 $n\geq2$ 时为\Term{bijection}{双射}。

再假设 $U$ 在 $X$ 中\Term{retrocompact}{逆紧}，则上述条件还\Term{equivalence}{等价}于：

(iii) 对 $Y$ 的每个\Term{geometric-point}{几何点} $\bar y$，\Term{canonical}{典范}\Term{morphism}{态射}
\[H^0(\overline X,\overline F)\longrightarrow H^0(\overline U,\overline F)\]
在 $n\geq1$ 时为\Term{injectivity}{单射}，在 $n\geq2$ 时为\Term{bijection}{双射}。

$2^\circ)$ 设 $F$ 为 $X$ 上的\Term{sheaf-of-groups}{群层}，$n\leq3$ 为整数，则下列条件\Term{equivalence}{等价}：

(i) \Term{canonical}{典范}\Term{morphism}{态射} $F\to i_*i^*F$ 在 $n\geq1$ 时为\Term{injectivity}{单射}，在 $n\geq2$ 时为\Term{bijection}{双射}；若 $n\geq3$，除上述条件外，\Term{pointed-sheaf}{带点集合层} $R^1i_*(i^*F)$ 还为零。

(ii) 对每个在 $X$ 上\Term{etale}{平展}的\Term{scheme}{概形} $X'$，\Term{canonical}{典范}\Term{morphism}{态射}
\[H^0(X',F')\longrightarrow H^0(U',F')\]
在 $n\geq1$ 时为\Term{injectivity}{单射}，在 $n\geq2$ 时为\Term{bijection}{双射}；此外，\Term{canonical}{典范}\Term{morphism}{态射}
\[H^1(X',F')\longrightarrow H^1(U',F')\]
在 $n\geq2$ 时为\Term{injectivity}{单射}，在 $n\geq3$ 时为\Term{bijection}{双射}。
\SourcePage{169}
(ii bis) 与(ii)相同，但 $n=2$ 时只假设 $H^0(X',F')\to H^0(U',F')$ 为\Term{bijection}{双射}。

再假设 $U$ 在 $X$ 中\Term{retrocompact}{逆紧}，则上述条件还\Term{equivalence}{等价}于：

(iii) 对 $Y$ 的每个\Term{geometric-point}{几何点} $\bar y$，\Term{canonical}{典范}\Term{morphism}{态射}
\[H^0(\overline X,\overline F)\longrightarrow H^0(\overline U,\overline F)\]
在 $n\geq1$ 时为\Term{injectivity}{单射}，在 $n\geq2$ 时为\Term{bijection}{双射}；若 $n\geq3$，除上述条件外，$H^1(\overline U,\overline F)$ 还为零。

$3^\circ)$ 设 $F$ 为次数有下界的\Term{abelian-sheaf}{阿贝尔层}\Term{complex}{复形}，$n$ 为整数，则下列条件\Term{equivalence}{等价}：

(i) 当 $p<n$ 时有 $\mathcal H_Y^p(F)=0$（参看1.0）。

(ii) 对每个在 $X$ 上\Term{etale}{平展}的\Term{scheme}{概形} $X'$，\Term{canonical}{典范}\Term{morphism}{态射}
\[H^p(X',F')\longrightarrow H^p(U',F')\]
在 $p<n-1$ 时为\Term{bijection}{双射}，在 $p=n-1$ 时为\Term{injectivity}{单射}。

假设 $U$ 在 $X$ 中\Term{retrocompact}{逆紧}，则上述条件还\Term{equivalence}{等价}于：

(iii) 对 $Y$ 的每个\Term{geometric-point}{几何点} $\bar y$，\Term{canonical}{典范}\Term{morphism}{态射}
\[H^p(\overline X,\overline F)\longrightarrow H^p(\overline U,\overline F)\]
在 $p<n-1$ 时为\Term{bijection}{双射}，在 $p=n-1$ 时为\Term{injectivity}{单射}。

若 $F$ 为\Term{abelian-sheaf}{阿贝尔层}且 $n\geq2$，则(i)、(ii)还\Term{equivalence}{等价}于：

(ii bis) 对每个在 $X$ 上\Term{etale}{平展}的\Term{scheme}{概形} $X'$，\Term{canonical}{典范}\Term{morphism}{态射} $H^p(X',F')\to H^p(U',F')$ 在 $p<n-1$ 时为\Term{bijection}{双射}。
\end{SGAresult}
\begin{proof}
$1^\circ)$ 显然(i)$\Leftrightarrow$(ii)。说明当 $U$ 在 $X$ 中\Term{retrocompact}{逆紧}时，(i)$\Leftrightarrow$(iii)。事实上，(i)\Term{equivalence}{等价}于：对 $X$ 的每个\Term{geometric-point}{几何点} $\bar y$，\Term{morphism}{态射} $F_{\bar y}\to(i_*i^*F)_{\bar y}$ 在 $n\leq1$ 时为\Term{injectivity}{单射}，在 $n\leq2$ 时为\Term{bijection}{双射}（SGA 4 VIII 3.6）。当 $\bar y$ 为 $U$ 的\Term{geometric-point}{几何点}时，此\Term{morphism}{态射}总是\Term{bijection}{双射}，故可只考虑 $Y$ 的\Term{geometric-point}{几何点} $\bar y$。由 $i$ 的\Term{quasi-compact}{拟紧}性及（SGA 4 VIII 5.3），\Term{morphism}{态射}
\[F_{\bar y}\longrightarrow(i_*i^*F)_{\bar y}\]
\Term{canonical}{典范}地等同于 $H^0(\overline X,\overline F)\to H^0(\overline U,\overline F)$，故(i)与(iii)\Term{equivalence}{等价}。
\SourcePage{170}
$2^\circ)$ (i)$\Rightarrow$(ii)。关于 $H^0$ 的断言由 $1^\circ)$ 得到。设 $i'$ 为 $U'$ 在 $X'$ 中的\Term{canonical}{典范}\Term{immersion}{浸入}；关于 $H^1$ 的断言由\Term{exact-sequence}{正合列}（SGA 4 XII 3.2）得到：
\[0\longrightarrow H^1(X',i'_*i'^*F')\longrightarrow H^1(U',F')\longrightarrow H^0(X',R^1i'_*(i'^*F')).\]
(ii bis)$\Rightarrow$(i)。由 $1^\circ)$，只须说明 $n\geq3$ 时 $R^1i_*(i^*F)=0$。而 $R^1i_*(i^*F)$ 是\Term{presheaf}{预层} $X'\mapsto H^1(U',F')$ 的\Term{associated-sheaf}{伴随层}，依假设也就是\Term{presheaf}{预层} $X'\mapsto H^1(X',F')$ 的\Term{associated-sheaf}{伴随层}，后者为零。

(i)$\Leftrightarrow$(iii)。考虑到 $1^\circ)$，只须说明 $R^1i_*(i^*F)=0$ \Term{equivalence}{等价}于对 $Y$ 的每个\Term{geometric-point}{几何点} $\bar y$ 有 $H^1(\overline U,\overline F)=0$。由于 $R^1i_*(i^*F)$ 在 $Y$ 外为零，其为零\Term{equivalence}{等价}于对 $Y$ 的每个\Term{geometric-point}{几何点} $\bar y$ 有 $(R^1i_*(i^*F))_{\bar y}=0$。只须注意到 $i$ \Term{quasi-compact}{拟紧}，故 $(R^1i_*(i^*F))_{\bar y}=H^1(\overline U,\overline F)$（SGA 4 VIII 5.3）。

$3^\circ)$ (i)$\Rightarrow$(ii)。设 $X'$ 为 $X$ 上\Term{etale}{平展}的\Term{scheme}{概形}；有\Term{exact-sequence}{正合列}（SGA 4 V 4.5）
\[\longrightarrow H^p_{Y'}(X',F')\longrightarrow H^p(X',F')\longrightarrow H^p(U',F')\longrightarrow.\tag{*}\]
故(ii)\Term{equivalence}{等价}于：对 $p<n$ 及每个在 $X$ 上\Term{etale}{平展}的\Term{scheme}{概形} $X'$，有 $H^p_{Y'}(X',F')=0$。考虑\Term{spectral-sequence}{谱序列}
\[E_2^{pq}=H^p(X',\mathcal H_Y^q(F))\Longrightarrow H^*_{Y'}(X',F').\]
由假设，$q<n$ 时 $\mathcal H_Y^q(F)=0$，故 $p+q<n$ 时 $E_2^{pq}=0$，从而 $p<n$ 时 $H^p_{Y'}(X',F')=0$。

(ii)$\Rightarrow$(i)。\Term{sheaf}{层} $\mathcal H_Y^p(F)$ 是\Term{presheaf}{预层} $X'\mapsto H^p_{Y'}(X',F')$ 的\Term{associated-sheaf}{伴随层}。已经指出(ii)\Term{equivalence}{等价}于对 $p<n$ 及所有在 $X$ 上\Term{etale}{平展}的\Term{scheme}{概形} $X'$ 有 $H^p_{Y'}(X',F')=0$，故确有 $\mathcal H_Y^p(F)=0$。

(i)$\Leftrightarrow$(iii)。\Term{sheaf}{层} $\mathcal H_Y^p(F)$ 集中在 $Y$ 上，故其为零\Term{equivalence}{等价}于它在 $Y$ 每个\Term{geometric-point}{几何点} $\bar y$ 的\Term{stalk}{茎} $(\mathcal H_Y^p(F))_{\bar y}$ 为零。由于 $i$ \Term{quasi-compact}{拟紧}，由（SGA 4 VIII 5.2）得 $(\mathcal H_Y^p(F))_{\bar y}=H^p_{\overline Y}(\overline X,\overline F)$。结合 $\overline X$ 上与\Term{exact-sequence}{正合列}(*)对应的\Term{exact-sequence}{正合列}，便得(i)与(iii)\Term{equivalence}{等价}。

当 $F$ 为\Term{abelian-sheaf}{阿贝尔层}时，(ii bis)$\Rightarrow$(ii)。只须再说明 $\mathcal H_Y^{n-1}(F)=0$。当 $n>2$ 时，这个\Term{sheaf}{层}是\Term{presheaf}{预层} $X'\mapsto H^{n-2}(U',F')=H^{n-2}(X',F')$ 的\Term{associated-sheaf}{伴随层}，故为零。$n=2$ 的情形来自 $\mathcal H_Y^1(F)$ 是\Term{morphism}{态射} $F\to i_*i^*F$ 的\Term{cokernel}{余核}这一事实。
\end{proof}
\begin{SGAresult}{定义}{1.2}
记号同1.1。若 $F$ 满足1.1中\Term{equivalence}{等价}条件(i)、(ii)，则称 $F$ 的 $Y$-\Term{etale-depth}{平展深度} $\geq n$，记为
\[\operatorname{prof}_Y(F)\geq n.\]
若 $F$ 为\Term{abelian-sheaf}{阿贝尔层}\Term{complex}{复形}，则称满足 $\operatorname{prof}_Y(F)\geq n$ 的整数 $n$ 的\Term{supremum}{上确界}为 $F$ 的 $Y$-\Term{etale-depth}{平展深度}。
\SourcePage{171}
若 $F$ 为\Term{sheaf-of-sets}{集合层}或未必交换的\Term{sheaf-of-groups}{群层}，也使用同样的记号（因此分别有 $0\leq\operatorname{prof}_Y(F)\leq2$ 或 $0\leq\operatorname{prof}_Y(F)\leq3$，只要上下文不致混淆此处所考虑的三种变体）。

若 $L$ 是素数\Term{set}{集合}，且对每个形如 $\mathbf Z/\ell\mathbf Z$、$\ell\in L$ 的\Term{constant-sheaf}{常值层}有 $\operatorname{prof}_Y(\mathbf Z/\ell\mathbf Z)\geq n$，则称 $X$ 关于 $L$ 的 $Y$-\Term{etale-depth}{平展深度} $\geq n$，记为
\[\operatorname{prof}_Y^L(X)\geq n.\]
据此显然可定义 $X$ 关于 $L$ 的 $Y$-\Term{etale-depth}{平展深度}。若 $L=\mathbf P$ 为全体素数的\Term{set}{集合}，且不会与（EGA IV 5.7.1）中的记号（关于 $F=\mathcal O_X$ 的情形）混淆，则略去记号中的 $L$；否则写作 $\operatorname{prof\text{ét}}(X)$。

最后，若对 $X$ 上的每个\Term{finite}{有限}常值 $L$-\Term{sheaf-of-groups}{群层} $F$ 都有 $\operatorname{prof}_Y(F)\geq3$，则称 $X$ 关于 $L$ 的 $Y$-\Term{homotopical-depth}{同伦深度} $\geq3$，记为
\[\operatorname{prof\,hop}_Y^L(X)\geq3.\]
若 $L=\mathbf P$，则略去记号中的 $L$。
\end{SGAresult}
\begin{SGAresult}{推论}{1.3}
在1.1的条件下，若 $\operatorname{prof}_Y(F)\geq n$，则对 $Y$ 的每个\Term{closed-subset}{闭子集} $Z$ 有 $\operatorname{prof}_Z(F)\geq n$。
\end{SGAresult}
\begin{proof}
例如在 $F$ 为次数有下界的\Term{abelian-sheaf}{阿贝尔层}\Term{complex}{复形}的情形下论证。使用1.1 $3^\circ)$(ii)。令 $V=X-Z$，对每个整数 $p$ 考虑\Term{morphism}{态射}列
\[H^p(X,F)\xrightarrow{f}H^p(V,F)\xrightarrow{g}H^p(U,F).\]
由假设，$g$ 与 $f\circ g$ 在 $p<n-1$ 时为\Term{bijection}{双射}，在 $p=n-1$ 时为\Term{injectivity}{单射}，故 $f$ 亦然。用 $X$ 上\Term{etale}{平展}的\Term{scheme}{概形} $X'$ 代替 $X$ 后论证仍适用，故1.3得证。
\end{proof}
\begin{SGAresult}{推论}{1.4}
记号同1.1 $2^\circ)$。若 $X'$ 为 $X$ 上的\Term{scheme}{概形}，以 $\Phi'$ 记把 $X'$ 的\Term{etale-covering}{平展覆盖}映为其在 $U'$ 上的限制的\Term{functor}{函子}，以 $\Phi'_{F'}$ 记把 $F'$ 下的\Term{torsor}{挠子}\footnote{即旧术语中的“\Term{principal-homogeneous-bundle}{主齐性纤维丛}”。}映为其在 $U'$ 上的限制的\Term{functor}{函子}。则下列条件\Term{equivalence}{等价}：

(i) $\operatorname{prof}_Y(F)\geq1$（分别 $\geq2$、$\geq3$）。

(ii) 对每个在 $X$ 上\Term{etale}{平展}的\Term{scheme}{概形} $X'$，\Term{functor}{函子} $\Phi'_{F'}$ 是\Term{faithful}{忠实}的（分别\Term{fully-faithful}{全忠实}、\Term{equivalence-of-categories}{范畴等价}）。

特别地，$\operatorname{prof}_Y(X)\geq1$（分别 $\operatorname{prof}_Y(X)\geq2$、$\operatorname{prof\,hop}_Y(X)\geq3$）的充要条件是 $\Phi'$ 为\Term{faithful}{忠实}\Term{functor}{函子}（分别\Term{fully-faithful}{全忠实}\Term{functor}{函子}、\Term{equivalence-of-categories}{范畴等价}）。
\end{SGAresult}
\SourcePage{172}
\begin{proof}
这由1.1 $2^\circ)$(ii)得到，只须利用 $H^1(X',F')$ 是 $F'$ 下的\Term{torsor}{挠子}的\Term{isomorphism-class}{同构类}\Term{set}{集合}这一解释（SGA 4 VII 2），以及\Term{scheme}{概形}的次数为 $n$ 的\Term{etale-covering}{平展覆盖} $Z$ 与群为\Term{symmetric-group}{对称群} $S_n$ 的\Term{galois-covering}{Galois 主覆盖}之间的对应：对应于 $Z$ 的\Term{covering}{覆盖}为 $\operatorname{Isom}_X(Z_0,Z)$，其中 $Z_0$ 是 $X$ 的次数为 $n$ 的\Term{trivial-covering}{平凡覆盖}。
\end{proof}
\begin{SGAresult}{推论}{1.5}
在1.1 $3^\circ)$ 的条件下，若 $\operatorname{prof}_Y(F)\geq n$，则有
\[H_Y^n(X,F)\simeq H^0(X,\mathcal H_Y^n(F)).\]
\end{SGAresult}
\begin{proof}
由\Term{spectral-sequence}{谱序列} $E_2^{pq}=H^p(X,\mathcal H_Y^q(F))\Longrightarrow H_Y^*(X,F)$ 得到。事实上，依假设 $q<n$ 时 $E_2^{pq}=0$，故
\[H_Y^n(X,F)=E_2^{0n}=H^0(X,\mathcal H_Y^n(F)).\]
\end{proof}
\textbf{注记1.6。}
a) 以1.1的\Term{equivalence}{等价}条件(i)、(ii)表述的 $Y$-\Term{depth}{深度}概念对任意\Term{site}{景}都有意义。特别地，若 $X$ 为配备 Zariski \Term{topology}{拓扑}的\Term{locally-noetherian}{局部诺特}\Term{scheme}{概形}，$F$ 为\Term{coherent-sheaf}{凝聚} $\mathcal O_X$-\Term{module-sheaf}{模层}，便得到通常的 $Y$-\Term{depth}{深度}概念，即 $Y$ 各点\Term{depth}{深度}的\Term{infimum}{下确界}（第 III 讲）。

b) $n\leq2$ 时，$X$ 的 $Y$-\Term{etale-depth}{平展深度}概念与 $L$ 无关。$n=1$ 时，它只表示 $U$ 在 $X$ 中\Term{dense}{稠密}。事实上，此条件是 $\operatorname{prof}_Y(F)\geq1$ 的必要条件，也是充分条件：可以假设 $X$ \Term{reduced}{既约}，此时 $U$ 在 $X$ 中\Term{dense}{稠密}的条件在\Term{etale}{平展}\Term{base-change}{基变换}下保持（EGA IV 11.10.5 (ii) b)）。若 $U$ 在 $X$ 中\Term{retrocompact}{逆紧}，则 $\operatorname{prof}_Y(X)\geq1$ 还\Term{equivalence}{等价}于 $Y$ 不含 $X$ 的任何\Term{maximal-point}{极大点}（EGA I 6.6.5）。$n=2$ 且 $U$ 在 $X$ 中\Term{retrocompact}{逆紧}时，条件 $\operatorname{prof}_Y(X)\geq2$ \Term{equivalence}{等价}于：对 $Y$ 的每个\Term{geometric-point}{几何点} $\bar y$，$\overline U$ 为\Term{nonempty}{非空}\Term{connected}{连通}的，即“对于\Term{etale-topology}{平展拓扑}，$Y$ \Term{local-scheme}{局部}不使 $X$ 不\Term{connected}{连通}”。

c) 若 $X$ 关于 $L$ 的 $Y$-\Term{depth}{深度} $\geq n$，且 $U$ 在 $X$ 中\Term{retrocompact}{逆紧}，则对 $X$ 上的每个\Term{locally-constant}{局部常值} $L$-\Term{torsion}{挠}\Term{abelian-sheaf}{阿贝尔层} $F$，有 $\operatorname{prof}_Y(X)\geq n$。事实上，性质 $\operatorname{prof}_Y(F)\geq n$ 对\Term{etale-topology}{平展拓扑}是\Term{local-scheme}{局部}的，故可设 $F$ 为\Term{constant-sheaf}{常值层}；于是 $F$ 是形如 $\mathbf Z/p^m\mathbf Z$（$m>0$、$p\in L$）的层的\Term{finite}{有限}\Term{direct-sum}{直和}所组成的\Term{filtered-colimit}{滤过归纳极限}。由1.1 (iii)及（SGA 4 VII 3.3），可归结到 $F=\mathbf Z/p^m\mathbf Z$，再对 $m$ 归纳，归结到 $F=\mathbf Z/p\mathbf Z$，此时断言由定义得到。

d) 由1.4，若 $\operatorname{prof}_Y(X)\geq3$，则偶对 $(X,Y)$ 在第 X 讲3.1的意义下是\Term{pure-pair}{纯偶对}。实际上，实践中遇到的\Term{pure-pair}{纯偶对}（参看第 X 讲3.4）满足更强的\Term{homotopical-depth}{同伦深度} $\geq3$ 条件，因此可用这一概念更\Term{effective}{有效}地代替\Term{pure-pair}{纯偶对}的概念。
\SourcePage{173}
e) 设 $F$ 为\Term{abelian-sheaf}{阿贝尔层}\Term{complex}{复形}，$T(F)$ 为对 $F$ 应用\Term{translation-functor}{平移函子}（[3]）所得的\Term{complex}{复形}，则显然有
\[\operatorname{prof}_Y(T(F))=\operatorname{prof}_Y(F)-1.\]
f) 指出 Artin--Mazur 最近的工作（[1]）允许对任意整数 $n$ 定义\Term{homotopical-depth}{同伦深度} $\geq n$ 的概念（不限于 $n\geq3$）。

g) 在1.1 $3^\circ)$ 的条件下，$\operatorname{prof}_Y(F)=\infty$ 的充要条件是在 $D^+(X)$ 中有 $F\xrightarrow{\sim}Ri_*(i^*F)$。事实上，$\mathcal H_Y^p(F)$ 是\Term{canonical}{典范}\Term{morphism}{态射} $F\to Ri_*i^*(F)$ 的锥（即映射柱）的\Term{cohomology-sheaf}{上同调层}。
\begin{SGAresult}{定义}{1.7}
设 $X$ 为\Term{scheme}{概形}，$x$ 为 $X$ 的点，$\bar x$ 为 $x$ 上方的\Term{geometric-point}{几何点}，$\overline X$ 为 $X$ 在 $\bar x$ 处的\Term{strict-localization}{严格局部化}。与前面一样，$F$ 表示 $X$ 上的\Term{sheaf-of-sets}{集合层}、\Term{sheaf-of-groups}{群层}，或次数有下界的\Term{abelian-sheaf}{阿贝尔层}\Term{complex}{复形}；$\overline F$ 为其在 $\overline X$ 上的\Term{inverse-image}{逆像}，$L$ 为素数\Term{set}{集合}。若 $\operatorname{prof}_{\bar x}(\overline F)>n$（分别 $\operatorname{prof}_{\bar x}^L(\overline X)>n$、$\operatorname{prof\,hop}_{\bar x}^L(\overline X)\geq3$），则称 $F$ 在点 $x$ 的\Term{etale-depth}{平展深度} $\geq n$（分别 $X$ 在 $x$ 关于 $L$ 的\Term{etale-depth}{平展深度} $\geq n$、关于 $L$ 的\Term{homotopical-depth}{同伦深度} $\geq3$），记为 $\operatorname{prof}_x(F)\geq n$（分别 $\operatorname{prof}_x^L(X)\geq n$、$\operatorname{prof\,hop}_x^L(X)\geq3$）。据此显然可定义整数 $\operatorname{prof}_x^L(X)$，以及在 $F$ 为\Term{abelian-sheaf}{阿贝尔层}\Term{complex}{复形}时的整数 $\operatorname{prof}_x(F)$。若 $L$ 为全体素数的\Term{set}{集合}，则略去 $\operatorname{prof}_x^L(X)$ 中的 $L$，除非可能与（EGA IV 5.7.1）中的记号混淆，此时写作 $\operatorname{prof\text{ét}}_x(X)$。
\end{SGAresult}
于是有如下\Term{depth}{深度}的逐点刻画。
\begin{SGAresult}{定理}{1.8}
设 $X$ 为\Term{scheme}{概形}，$Y$ 为其\Term{closed-subset}{闭子集}，且 $U=X-Y$ 在 $X$ 中\Term{retrocompact}{逆紧}。若 $F$ 是 $X$ 上的\Term{sheaf-of-sets}{集合层}、\Term{sheaf-of-groups}{群层}，或次数有下界的\Term{abelian-sheaf}{阿贝尔层}\Term{complex}{复形}，则
\[\operatorname{prof}_Y(F)=\inf_{y\in Y}\operatorname{prof}_y(F).\]
\end{SGAresult}
\begin{proof}
\textbf{1.8.1。}先说明对 $Y$ 的每点 $y$ 有 $\operatorname{prof}_y(F)\geq\operatorname{prof}_Y(F)$。设 $\bar y$ 为 $y$ 上方的\Term{geometric-point}{几何点}，$\overline X$ 为 $X$ 在 $\bar y$ 处的\Term{strict-localization}{严格局部化}，$\overline F$ 与 $\overline Y$ 为 $F$ 与 $Y$ 在 $\overline X$ 上的\Term{inverse-image}{逆像}。由1.7及1.3有
\[\operatorname{prof}_y(F)=\operatorname{prof}_{\bar y}(\overline F)\geq\operatorname{prof}_{\overline Y}(\overline F)\geq\operatorname{prof}_Y(F),\]
最后一个不等式通过1.1中的条件(iii)及\Term{strict-localization}{严格局部化}构造的传递性使用了“$U$ \Term{retrocompact}{逆紧}”这一假设。
\SourcePage{174}
\textbf{1.8.2。}反之，设对 $Y$ 的每点 $y$ 有 $\operatorname{prof}_y(F)\geq n$（$n$ 为整数），说明 $\operatorname{prof}_Y(F)\geq n$。

先回顾如下熟知结果（SGA 4 VIII）：
\end{proof}
\begin{SGAresult}{引理}{1.8.2.1}
设 $X$ 为\Term{scheme}{概形}，$F$ 为 $X$ 上的\Term{sheaf-of-sets}{集合层}（分别 $G\to F$ 为 $X$ 上\Term{sheaf-of-sets}{集合层}的\Term{monomorphism}{单态射}）。$F$ 的两个\Term{section}{截面} $s,s'$ 相等（分别 $F$ 的\Term{section}{截面} $s$ 来自 $G$ 的\Term{section}{截面}）的充要条件是\Term{local-scheme}{局部}成立。特别地，若 $s,s'$ 为 $F$ 的两个\Term{section}{截面}，则存在 $X$ 的最大\Term{open-subset}{开集} $V$，使二者在 $V$ 上相等（分别若 $s$ 为 $F$ 在 $X$ 上的\Term{section}{截面}，则存在 $X$ 的最大\Term{open-subset}{开集} $V$，使 $s|_V$ 来自 $G$ 在 $V$ 上的\Term{section}{截面}）。此\Term{open-subset}{开集}也就是这样的点 $x\in X$ 的\Term{set}{集合}：取 $x$ 上方的\Term{geometric-point}{几何点} $\bar x$，$s,s'$ 在\Term{stalk}{茎} $F_{\bar x}$ 中的像相等（分别 $s$ 在 $F_{\bar x}$ 中的像来自 $G_{\bar x}$ 的元素）。
\end{SGAresult}
\begin{proof}[定理1.8的证明（续）]
回到1.8的证明。

$1^\circ)$ $F$ 为\Term{sheaf-of-sets}{集合层}的情形。若 $n=1$，只须说明\Term{canonical}{典范}\Term{morphism}{态射} $H^0(X,F)\to H^0(U,F)$ 为\Term{injectivity}{单射}，因为用 $X$ 上\Term{etale}{平展}的\Term{scheme}{概形}代替 $X$ 后结论仍适用。设 $s,s'$ 为 $F$ 在 $X$ 上的两个\Term{section}{截面}，在 $H^0(U,F)$ 中具有相同的像；以 $V$ 记使二者相等的最大\Term{open-subset}{开集}，显然 $V\supset U$。假设 $V\ne X$，取 $X-V$ 的一个\Term{maximal-point}{极大点} $y$，$\bar y$ 为其上方的\Term{geometric-point}{几何点}，$\overline X$ 为 $X$ 在 $\bar y$ 处的\Term{strict-localization}{严格局部化}，$\overline V,\overline F$ 为 $V,F$ 在 $\overline X$ 上的\Term{inverse-image}{逆像}。由 $y$ 的选择，$\overline X-\bar y=\overline V$，而依假设
\[H^0(\overline X,\overline F)\longrightarrow H^0(\overline X-\bar y,\overline F)=H^0(\overline V,\overline F)\]
为\Term{injectivity}{单射}。故 $s,s'$ 在 $y$ 处相等，矛盾。若 $n=2$，由前述结论，只须说明 $H^0(X,F)\to H^0(U,F)=H^0(X,i_*i^*F)$ 为\Term{surjectivity}{满射}，其中 $i$ 为 $U$ 在 $X$ 中的\Term{canonical}{典范}\Term{immersion}{浸入}。设 $s$ 为 $i_*i^*F$ 在 $X$ 上的\Term{section}{截面}，$V$ 为使其来自 $F$ 的\Term{section}{截面}的最大\Term{open-subset}{开集}。假设 $V\ne X$，取 $X-V$ 的\Term{maximal-point}{极大点} $y$。沿用上述记号，由假设，\Term{canonical}{典范}\Term{morphism}{态射}
\[H^0(\overline X,\overline F)\longrightarrow H^0(\overline X-\bar y,\overline F)=H^0(\overline V,\overline F)\]
为\Term{bijection}{双射}，故 $s|_{\overline V}$ 可延拓到 $y$，矛盾。这就证明了 $1^\circ)$。

$2^\circ)$ $F$ 为\Term{sheaf-of-groups}{群层}的情形。由 $1^\circ)$，只须说明 $n=3$ 时
\[H^1(X,F)=H^1(X,i_*i^*F)\longrightarrow H^1(U,F)\]
\SourcePage{175}
为\Term{bijection}{双射}。由 $1^\circ)$ 及1.1 $2^\circ)$(ii bis)，已知它为\Term{injectivity}{单射}。为证\Term{surjectivity}{满射}，使用\Term{exact-sequence}{正合列}（SGA 4 XII 3.2）
\[0\longrightarrow H^1(X,i_*i^*F)\longrightarrow H^1(U,F)\xrightarrow{d}H^0(X,R^1i_*(i^*F)).\]
设 $s\in H^1(U,F)$，$V\supset U$ 为使 $d(s)=0$ 的最大\Term{open-subset}{开集}；它也是使 $s$ 来自 $H^1(V,F)$ 的元素的最大\Term{open-subset}{开集}。假设 $V\ne X$，取 $X-V$ 的\Term{maximal-point}{极大点} $y$。若 $\overline X$ 为 $X$ 在 $y$ 上方的\Term{geometric-point}{几何点} $\bar y$ 处的\Term{strict-localization}{严格局部化}，则按显然的记号有\Term{exact-sequence}{正合列}
\[0\longrightarrow H^1(\overline X,\bar i_*(\bar i^*\overline F))\longrightarrow H^1(\overline U,\overline F)\xrightarrow{\bar d}H^0(\overline X,R^1\bar i_*(\bar i^*\overline F)).\]
因 $i:\overline U\to\overline X$ \Term{quasi-compact}{拟紧}，$R^1\bar i_*(\bar i^*\overline F)$ 是 $R^1i_*(i^*F)$ 沿 $\overline X\to X$ 的\Term{inverse-image}{逆像}，故
\[H^0(\overline X,R^1\bar i_*(\bar i^*\overline F))=(R^1i_*(i^*F))_{\bar y}.\]
依假设及 $y\in Y$，\Term{morphism}{态射} $H^1(\overline X,\overline F)\to H^1(\overline V,\overline F)$ 为\Term{bijection}{双射}。$s$ 在 $H^1(\overline U,\overline F)$ 中的像 $\bar s$ 按 $V$ 的定义可延拓到 $\overline V$，故也可延拓到 $\overline X$。于是 $\bar d(\bar s)=0$，即 $d(s)$ 在\Term{geometric-stalk}{几何茎} $(R^1i_*(i^*F))_{\bar y}$ 中的像为零，与 $V$ 的定义矛盾。$2^\circ)$ 得证。

$3^\circ)$ $F$ 为次数有下界的\Term{abelian-sheaf}{阿贝尔层}\Term{complex}{复形}的情形。对 $n$ 归纳。由于 $F$ 的次数有下界，$n$ 充分小时结论成立。现设 $\operatorname{prof}_Y(F)\geq n-1$，并已知对 $Y$ 的每点 $y$ 有 $\operatorname{prof}_y(F)\geq n$，来说明 $\operatorname{prof}_Y(F)\geq n$。只须证明\Term{canonical}{典范}\Term{morphism}{态射}
\[H^{n-2}(X,F)\longrightarrow H^{n-2}(U,F)\tag{*}\]
为\Term{surjectivity}{满射}，且
\[H^{n-1}(X,F)\longrightarrow H^{n-1}(U,F)\tag{**}\]
为\Term{injectivity}{单射}（因为用 $X$ 上\Term{etale}{平展}的\Term{scheme}{概形}代替 $X$ 后结论仍适用）。

a) (*)的\Term{surjectivity}{满射}性。证明类似于 $2^\circ)$。由1.5及（SGA 4 V 4.5），有\Term{exact-sequence}{正合列}
\[H^{n-2}(X,F)\longrightarrow H^{n-2}(U,F)\xrightarrow d H_Y^{n-1}(X,F)=H^0(X,\mathcal H_Y^{n-1}(F)).\]
设 $s\in H^{n-2}(U,F)$，$V\supset U$ 为使 $d(s)=0$ 的最大\Term{open-subset}{开集}，它也是使 $s$ 能延拓到 $H^{n-2}(V,F)$ 的最大\Term{open-subset}{开集}。假设 $V\ne X$，取 $X-V$ 的\Term{maximal-point}{极大点} $y$，$\overline X$ 为 $X$ 在 $y$ 上方的\Term{geometric-point}{几何点} $\bar y$ 处的\Term{strict-localization}{严格局部化}。因为 $i:U\to X$ \Term{quasi-compact}{拟紧}，
\SourcePage{176}
$\mathcal H_Y^{n-1}(F)$ 的构造与\Term{base-change}{基变换} $\overline X\to X$ 交换，故有\Term{exact-sequence}{正合列}（采用显然的记号）
\[H^{n-2}(\overline X,\overline F)\longrightarrow H^{n-2}(\overline U,\overline F)\xrightarrow{\bar d}H_{\overline Y}^{n-1}(\overline X,\overline F)=(\mathcal H_Y^{n-1}(F))_{\bar y},\]
最后一个等式由 $U$ 的\Term{retrocompact}{逆紧}性假设得到。依假设有\Term{isomorphism}{同构}
\[H^{n-2}(\overline X,\overline F)\xrightarrow\sim H^{n-2}(\overline X-\bar y,\overline F)=H^{n-2}(\overline V,\overline F).\]
故 $s$ 在 $H^{n-2}(\overline U,\overline F)$ 中的像 $\bar s$ 按 $V$ 的定义可延拓到 $H^{n-2}(\overline V,\overline F)$，因而也可延拓到 $H^{n-2}(\overline X,\overline F)$。这说明 $\bar d(\bar s)=0$，即 $d(s)$ 在 $\bar y$ 处为零，矛盾。

b) (**)的\Term{injectivity}{单射}性。由(*)的\Term{surjectivity}{满射}性得到\Term{exact-sequence}{正合列}
\[0\longrightarrow H^0(X,\mathcal H_Y^{n-1}(F))\longrightarrow H^{n-1}(X,F)\longrightarrow H^{n-1}(U,F),\]
须说明每个元素 $s\in H^0(X,\mathcal H_Y^{n-1}(F))$ 为零。令 $V$ 为使 $s=0$ 的最大\Term{open-subset}{开集}。假设 $V\ne X$，取 $X-V$ 的\Term{maximal-point}{极大点} $y$，$\overline X$ 为 $X$ 在 $y$ 上方的\Term{geometric-point}{几何点} $\bar y$ 处的\Term{strict-localization}{严格局部化}。由归纳假设及1.8.1，有 $\operatorname{prof}_{\overline Y}(\overline F)\geq\operatorname{prof}_Y(F)\geq n-1$，故下列图中的映射 $e$ 为\Term{injectivity}{单射}：
\[H^0(\overline X,\mathcal H_{\overline Y}^{n-1}(\overline F))=(\mathcal H_Y^{n-1}(F))_{\bar y}\xrightarrow e H^{n-1}(\overline X,\overline F)\xrightarrow f H^{n-1}(\overline V,\overline F).\]
由假设，$f$ 也为\Term{injectivity}{单射}；左边的等式由 $U$ 的\Term{retrocompact}{逆紧}性假设得到。以 $\bar s$ 记 $s$ 在 $(\mathcal H_Y^{n-1}(F))_{\bar y}$ 中的像。由于 $s$ 在 $V$ 上消失，$f\cdot e(\bar s)=0$，故 $\bar s=0$，这与 $y$ 的选择矛盾。证明完成。
\end{proof}
\textbf{注记1.9。}若把\Term{scheme}{概形} $X$ 的\Term{etale-topos}{平展拓扑斯}换成“\Term{locally-finite-type-topos}{局部有限型拓扑斯}”，即能由\Term{locally-finite-type-site}{局部有限型景}（SGA 4 VI 1.1）定义的\Term{topos}{拓扑斯}，则类似1.8的结果无疑仍成立。为此须使用 P. Deligne 的一个结果（SGA 4 VI.9），即这种\Term{topos}{拓扑斯}中有“足够多的\Term{fiber-functor}{纤维函子}”。

下面由1.8推出若干可以确定\Term{etale-depth}{平展深度}的重要情形。
\SourcePage{177}
\begin{SGAresult}{定理}{1.10}
\textbf{（\Term{cohomology}{上同调}半\Term{purity-theorem}{纯性定理}）}
\footnote{编者注：Gabber 后来于1994年证明了 Grothendieck 的\Term{absolute-cohomological-purity}{绝对上同调纯性}猜想：若 $Y$ 为绝对\Term{noetherian-scheme}{诺特概形}的纯余维 $c$ 的\Term{closed-subscheme}{闭子概形}，$n$ 是在 $X$ 上\Term{invertible}{可逆}的整数，置 $\Lambda=\mathbf Z/n\mathbf Z$，则 $q\ne2c$ 时 $\mathcal H_Y^q(\Lambda)=0$，而 $q=2c$ 时它等于 $\Lambda_Y(-c)$（\Term{tate-twist}{Tate 扭转}）。参看 Fujiwara K., ``A Proof of the Absolute Purity Conjecture (after Gabber)'', \emph{Algebraic geometry 2000, Azumino (Hotaka)}, Adv. Stud. in Pure Math., vol. 36, 2002, pp. 153--183。关于对\Term{dualizing-complex}{对偶化复形}存在性的应用，参看 SGA 5, Lect. Notes in Math., vol. 589, Springer-Verlag, 1977, p. 1672，第 I 讲，以及上述文献第8节。对于 $n=\ell^\nu$ 且 $\ell$ 为在 $X$ 上\Term{invertible}{可逆}的充分大的素数的情形，此猜想早已通过关键地使用 $K$-理论得到证明（Thomason R. W., ``Absolute cohomological purity'', Bull. Soc. Math. France 112 (1984), no. 3, pp. 397--406）。Gabber 的证明也通过联系\Term{etale-cohomology}{平展上同调}和 $K$-理论的 Atiyah--Hirzebruch--Thomason \Term{spectral-sequence}{谱序列}使用了 $K$-理论，这一方法已用于 Thomason 的途径。除此之外，另一项基本论据是原书第181页注(5)所引 Lefschetz 定理的推广。}
以 $X$ 表示域 $k$ 上的\Term{smooth}{光滑}\Term{scheme}{概形}，或\Term{characteristic}{特征}零的\Term{excellent}{优秀}\Term{regular}{正则}\Term{scheme}{概形}（EGA IV 7.8.2）。（注意：若承认（SGA 4 XIX）意义下的\Term{resolution-of-singularities}{奇点消解}，则更一般地，只须假设 $X$ 为\Term{equal-characteristic}{等特征}的\Term{excellent}{优秀}\Term{regular}{正则}\Term{scheme}{概形}。）设 $Y$ 为 $X$ 的\Term{closed-subset}{闭子集}，$L$ 为不同于 $X$ 的\Term{characteristic}{特征}的素数的\Term{set}{集合}，则
\[\operatorname{prof}_Y^L(X)=2\operatorname{codim}(Y,X).\]
\end{SGAresult}
\begin{proof}
由1.8，有 $\operatorname{prof}_Y^L(X)=\inf_{y\in Y}\operatorname{prof}_y(X)$。另一方面，$\operatorname{codim}(Y,X)=\inf_{y\in Y}\dim\mathcal O_{X,y}$，故归结为说明
\[\operatorname{prof}_y^L(X)=2\dim\mathcal O_{X,y},\]
这由（SGA 4 XVI 3.7 与 XIX 3.2）得到。
\end{proof}
\begin{SGAresult}{定理}{1.11}
\textbf{（\Term{homotopy}{同伦}\Term{purity-theorem}{纯性定理}）}
\footnote{编者注：最近 de Jong 与 Oort 得到了如下\Term{purity}{纯性}陈述：设 $\widetilde S\to S$ 为二维\Term{normal}{正规}\Term{noetherian-local-ring}{诺特局部环}的谱 $S$ 的\Term{resolution-singularities}{奇点消解}，$U$ 为 $S$ 中\Term{closed-point}{闭点} $s$ 的补集。再假设 $k(s)$ 为\Term{algebraically-closed-field}{代数闭域}。则对每个素数 $p$，特别当 $S$ 的\Term{characteristic}{特征}为 $p$ 时，限制\Term{morphism}{态射} $H^1_{\text{ét}}(\widetilde S,\mathbf Q_p)\to H^1_{\text{ét}}(U,\mathbf Q_p)$ 为\Term{bijection}{双射}（de Jong A. J. \& Oort F., ``Purity of the stratification by Newton polygons'', J. Amer. Math. Soc. 13 (2000), no. 1, pp. 209--241，定理3.2）。若 $k=\mathbf C$，$A$ 为曲面奇点的\Term{completion}{完备化}，此结果归于 Mumford（参看原书第158页[5]）。}
若 $X$ 为\Term{locally-noetherian}{局部诺特}\Term{scheme}{概形}，它是\Term{regular}{正则}的（分别，其\Term{local-ring}{局部环}为\Term{complete-intersection}{完全交}），$Y$ 为 $X$ 的\Term{closed-subset}{闭子集}，且 $\operatorname{codim}(Y,X)\geq2$（分别 $\geq3$），则
\[\operatorname{prof\,hop}_Y(X)\geq3.\]
\end{SGAresult}
\begin{proof}
由1.8，$\operatorname{prof\,hop}_Y(X)=\inf_{y\in Y}\operatorname{prof\,hop}_y(X)$。$X$ 在 $Y$ 各点的\Term{strictly-local-ring}{严格局部环}是\Term{dimension}{维数} $\geq2$ 的\Term{regular}{正则}环（分别\Term{dimension}{维数} $\geq3$ 的\Term{complete-intersection}{完全交}环）。由第 X 讲3.4的\Term{purity-theorem}{纯性定理}，$\operatorname{prof\,hop}_y(X)\geq3$，定理得证。
\end{proof}
\textbf{例1.12。}设 $X$ 为\Term{locally-noetherian}{局部诺特}\Term{scheme}{概形}，$Y$ 为其\Term{closed-subset}{闭子集}，$n=1$ 或2。若 $\operatorname{prof}_Y(\mathcal O_X)\geq n$（此处指\Term{coherent-sheaf}{凝聚层}意义的 $Y$-\Term{depth}{深度}，参看1.6 a)），则也有 $\operatorname{prof}_Y(X)\geq n$。$n=1$ 时显然，$n=2$ 时正是 Hartshorne 定理（第 III 讲1）。但 $n\geq3$ 时对应断言不成立。例如取\Term{characteristic}{特征}不为2的域上\Term{dimension}{维数} $\geq3$ 的\Term{affine-space}{仿射空间}，
\SourcePage{178}
使群 $\mathbf Z/2\mathbf Z$ 以关于原点的对称\Term{action}{作用}。令 $X$ 为其商，$Y=\{x\}$ 为原点在 $X$ 中的像。则 $\mathcal O_{X,x}$ 为\Term{cohen-macaulay-ring}{Cohen--Macaulay 环}，故 $\operatorname{prof}_x(\mathcal O_X)\geq3$；但去掉原点的\Term{affine-space}{仿射空间}是 $X-\{x\}$ 的\Term{etale-covering}{平展覆盖}，不能延拓为 $X$ 的\Term{etale-covering}{平展覆盖}，故由1.4，有 $\operatorname{prof}_Y(X)=2$。

下面的定理与（EGA IV 6.3.1）相对应。
\begin{SGAresult}{定理}{1.13}
设 $f:X\to S$ 为\Term{scheme}{概形}的\Term{morphism}{态射}，$Y$ 为 $X$ 的\Term{closed-subset}{闭子集}，$Z$ 为 $S$ 的\Term{closed-subset}{闭子集}，且 $f(Y)\subset Z$。假设 $X$ 在 $Y$ 各点的\Term{local-ring}{局部环}是\Term{noetherian}{诺特}的，且\Term{open-subset}{开集} $X-Y$ 与 $S-Z$ 分别在 $X$ 与 $S$ 中\Term{retrocompact}{逆紧}。设 $p,q,r$ 为整数，$p\geq-r$、$q\geq0$，$L$ 为素数\Term{set}{集合}，$F$ 为 $S$ 上的 $L$-\Term{torsion}{挠}\Term{abelian-sheaf}{阿贝尔层}\Term{complex}{复形}，其\Term{cohomology-sheaf}{上同调层} $\mathcal H^i(F)$ 在 $i<-r$ 时为零。假设：

a) $f$ 关于 $L$ \Term{locally-n-acyclic}{局部 $(p+q+r-2)$-无环}（SGA 4 XV 1.11）。

b) $\operatorname{prof}_Z(F)\geq p$。

c) 对 $Z$ 的每点 $s$，有 $\operatorname{prof}_{Y_s}^L(X_s)\geq q$。

则 $\operatorname{prof}_Y(f^*F)\geq p+q$。
\end{SGAresult}
需要如下引理。
\begin{SGAresult}{引理}{1.13.1}
设 $L$ 为素数\Term{set}{集合}，$n,r$ 为整数，$f:X\to S$ 为关于 $L$ \Term{locally-n-acyclic}{局部 $n$-无环}的\Term{morphism}{态射}。设 $F$ 为\Term{abelian-sheaf}{阿贝尔层}\Term{complex}{复形}，其\Term{cohomology-sheaf}{上同调层}为 $L$-\Term{torsion}{挠}，且 $i<-r$ 时 $\mathcal H^i(F)=0$；$Z$ 为 $S$ 的\Term{closed-subset}{闭子集}，$S-Z$ 在 $S$ 中\Term{retrocompact}{逆紧}，$T=f^{-1}(Z)$。则\Term{canonical}{典范}\Term{morphism}{态射}
\[f^*(\mathcal H_Z^i(F))\longrightarrow\mathcal H_T^i(f^*F)\]
在 $i<n-r+2$ 时为\Term{bijection}{双射}，在 $i=n-r+2$ 时为\Term{injectivity}{单射}。
\end{SGAresult}
\begin{proof}
置 $U=S-Z$、$V=X-T$，有\Term{cartesian-square}{笛卡尔方块}
\[\begin{tikzcd}V\arrow[r,"g"]\arrow[d,"k"']&U\arrow[d,"j"]\\X\arrow[r,"f"']&S.\end{tikzcd}\]
\SourcePage{179}
考虑如下各行\Term{exact}{正合}的\Term{commutative-diagram}{交换图}：
\[\begin{tikzcd}[column sep=small]
{}\arrow[r]&f^*(\mathcal H_Z^i(F))\arrow[r]\arrow[d]&f^*(\mathcal H^i(F))\arrow[r]\arrow[d,"\sim"]&f^*(\mathcal H^i(Rj_*(j^*F)))\arrow[r]\arrow[d]&{}\\
{}\arrow[r]&\mathcal H_T^i(f^*F)\arrow[r]&\mathcal H^i(f^*F)\arrow[r]&\mathcal H^i(Rk_*(k^*f^*F))\arrow[r]&{}.
\end{tikzcd}\]
由此归结为说明\Term{morphism}{态射}
\[f^*(\mathcal H^i(Rj_*(j^*F)))\longrightarrow\mathcal H^i(Rk_*(k^*f^*F))\]
在 $i<n-r+1$ 时为\Term{bijection}{双射}，在 $i=n-r+1$ 时为\Term{injectivity}{单射}。此\Term{morphism}{态射}来自下列\Term{hypercohomology}{超上同调}\Term{spectral-sequence}{谱序列}的\Term{morphism}{态射}：
\[\begin{tikzcd}[column sep=large]
f^*E_2^{p,q}=f^*(R^pj_*(\mathcal H^q(j^*F)))\arrow[r,Rightarrow]\arrow[d]&f^*(\mathcal H^*(Rj_*(j^*F)))\arrow[d]\\
E_2'{}^{p,q}=R^pk_*(\mathcal H^q(k^*f^*F))\arrow[r,Rightarrow]&\mathcal H^*(Rk_*(k^*f^*F)).
\end{tikzcd}\]
因 $j$ \Term{quasi-compact}{拟紧}，由（SGA 4 XV 1.10），\Term{morphism}{态射} $f^*(E_2^{p,q})\to E_2'{}^{p,q}$ 在 $p\leq n$ 时为\Term{bijection}{双射}，在 $p=n+1$ 时为\Term{injectivity}{单射}；特别地，在 $p+q\leq n-r$ 时为\Term{bijection}{双射}，在 $p+q=n-r+1$ 时为\Term{injectivity}{单射}。结论随即得到。
\end{proof}
\begin{proof}[定理1.13的证明]
回到1.13。令 $T=f^{-1}(Z)$。由1.13.1及条件a)，\Term{canonical}{典范}\Term{morphism}{态射} $f^*(\mathcal H_Z^i(F))\to\mathcal H_T^i(f^*F)$ 在 $i\leq p+q$ 时为\Term{isomorphism}{同构}。故由b)，$i<p$ 时 $\mathcal H_T^i(f^*F)=0$，而 $i<p+q$ 时，$\mathcal H_T^i(f^*F)$ 在 $T$ 上的限制是 $Z$ 上某\Term{sheaf}{层} $G^i$ 的\Term{inverse-image}{逆像}。设 $f_T:T\to Z$ 为 $f$ 在 $T$ 上的限制。由c)及下面的推论，$j<q$ 时 $\mathcal H_Y^j(f_T^*(G^i))=0$，从而
\[\mathcal H_Y^j(\mathcal H_T^i(f^*F))=0\qquad(i+j<p+q).\]
因为 $i+j<p+q$ 蕴涵或 $i<p$，于是 $\mathcal H_T^i(f^*F)=0$；或 $j<q$，于是 $\mathcal H_Y^j(\mathcal H_T^i(f^*F))=0$。用1.0的记号，有 $\underline\Gamma_Y=\underline\Gamma_Y\cdot\underline\Gamma_T$，故有\Term{spectral-sequence}{谱序列}
\[E_2^{i,j}=\mathcal H_Y^j(\mathcal H_T^i(f^*F))\Longrightarrow\mathcal H_Y^*(f^*F).\tag{1.13.2}\]
由于 $i+j<p+q$ 时 $E_2^{i,j}=0$，可知 $k<p+q$ 时 $\mathcal H_Y^k(f^*F)=0$。故只须证明如下推论，便证明了定理（此推论是1.13中令 $Z=S$、$r=p=0$ 且 $F$ 集中在次数0的特例）。
\end{proof}
\begin{SGAresult}{推论}{1.14}
设 $f:X\to S$ 为\Term{morphism}{态射}，$Y$ 为\Term{closed-subset}{闭子集}，补\Term{open-subset}{开集} $X-Y$ 在 $X$ 中\Term{retrocompact}{逆紧}，且 $X$ 在 $Y$ 各点的\Term{local-ring}{局部环}为\Term{noetherian}{诺特}环。设 $L$ 为素数\Term{set}{集合}，$q$ 为整数，$F$ 为 $S$ 上 $L$-\Term{torsion}{挠}\Term{abelian-sheaf}{阿贝尔层}。
\SourcePage{180}
假设 $f$ 关于 $L$ \Term{locally-n-acyclic}{局部 $(q-2)$-无环}，且对 $S$ 的每点 $s$ 有 $\operatorname{prof}_{Y_s}^L(X_s)\geq q$，则 $\operatorname{prof}_Y(f^*F)\geq q$。
\end{SGAresult}
\begin{proof}
$1^\circ)$ 归结到 $X,S$ 为\Term{strictly-local-scheme}{严格局部概形}、$f$ 为\Term{local-morphism}{局部态射}，且 $Y$ 为 $X$ 的一个\Term{closed-point}{闭点}的情形。

由1.8，为证1.14，须说明对 $Y$ 的每点 $y$，$\operatorname{prof}_y(f^*F)\geq q$。令 $s=f(y)$，$\bar s$ 为 $s$ 上方的\Term{geometric-point}{几何点}，$\bar y$ 为 $y$ 及 $\bar s$ 上方的\Term{geometric-point}{几何点}，$\overline X,\overline S$ 分别为 $X,S$ 在 $\bar y,\bar s$ 处的\Term{strict-localization}{严格局部化}，$\bar f:\overline X\to\overline S$ 为\Term{canonical}{典范}\Term{morphism}{态射}，$\overline F$ 为 $F$ 在 $\overline S$ 上的\Term{inverse-image}{逆像}。因 $\operatorname{prof}_y(f^*F)=\operatorname{prof}_{\bar y}(\bar f^*\overline F)$，只须说明用 $\bar f$（分别 $\{\bar y\}$、$\overline F$）代替 $f$（分别 $Y$、$F$）后，1.14的假设保持成立。\Term{retrocompact}{逆紧}性条件由 $\mathcal O_{X,x}$ 的\Term{noetherian}{诺特}性假设得到，后者蕴涵 $\overline X$ 为\Term{noetherian}{诺特}的。由（SGA 4 XV 1.10 (i)），$\bar f$ 仍关于 $L$ \Term{locally-n-acyclic}{局部 $(q-2)$-无环}。另一方面，$\overline X$ 在 $\bar s$ 上的\Term{fiber}{纤维} $(\overline X)_{\bar s}$ 等同于 $X_s$ 在 $\bar y$ 处的\Term{strict-localization}{严格局部化}，故满足 $\operatorname{prof}_{\bar y}((\overline X)_{\bar s})\geq q$。除闭\Term{fiber}{纤维}之外，$\overline S$-\Term{scheme}{概形} $\overline X$ 的其余\Term{fiber}{纤维}显然满足类似关系，归结完成。

$2^\circ)$ $X,S$ 为\Term{strictly-local-scheme}{严格局部概形}，$f$ 为\Term{local-homomorphism}{局部同态}，$Y$ 为 $X$ 的\Term{closed-point}{闭点}。设
\[g:Y=X-\{y\}\longrightarrow S\]
为 $U$ 的\Term{structural-morphism}{结构态射}。须说明\Term{canonical}{典范}\Term{morphism}{态射} $u_i:H^i(X,f^*F)\to H^i(U,f^*F)$ 在 $i\leq q-2$ 时为\Term{bijection}{双射}，在 $i=q-1$ 时为\Term{injectivity}{单射}。考虑\Term{commutative-diagram}{交换图}
\[\begin{tikzcd}H^i(X,f^*F)\arrow[rr,"u_i"]&&H^i(U,f^*F)\\&H^i(S,F)\arrow[ul,"v_i"]\arrow[ur,"w_i"']&\end{tikzcd}\]
\Term{morphism}{态射} $v_i$ 对所有 $i$ 显然为\Term{bijection}{双射}。另一方面，$g$ 关于 $L$ \Term{locally-n-acyclic}{局部 $(q-2)$-无环}；它的\Term{fiber}{纤维}也关于 $L$ 为 $(q-2)$-\Term{acyclic}{无环}，这是因为 $\operatorname{prof}_y(X_s)\geq q$，且 $f$ 的\Term{fiber}{纤维}关于 $L$ 为 $(q-2)$-\Term{acyclic}{无环}。由于 $X$ \Term{noetherian}{诺特}，$g$ \Term{quasi-compact}{拟紧}，由（SGA 4 XV 1.16），$g$ 关于 $L$ 为 $(q-2)$-\Term{acyclic}{无环}。故 $w_i$，从而 $u_i$，在 $i\leq q-2$ 时为\Term{bijection}{双射}，在 $i=q-1$ 时为\Term{injectivity}{单射}。1.14得证。
\end{proof}
\begin{SGAresult}{推论}{1.15}
设 $f:X\to S$ 为\Term{scheme}{概形}的\Term{morphism}{态射}，$L$ 为素数\Term{set}{集合}，$m,r$ 为整数，$F$ 为 $S$ 上 $L$-\Term{torsion}{挠}\Term{abelian-sheaf}{阿贝尔层}\Term{complex}{复形}，且 $i<-r$ 时 $\mathcal H^i(F)=0$。设 $x\in X$，$s=f(x)$，
\SourcePage{181}
并假设\Term{local-ring}{局部环} $\mathcal O_{X,x}$ 为\Term{noetherian}{诺特}环。若 $f$ 关于 $L$ \Term{locally-n-acyclic}{局部 $m$-无环}，则有
\[\operatorname{prof}_x(f^*F)\geq\inf\bigl(\operatorname{prof}_s(F)+\operatorname{prof}_x^L(X_s),n\bigr),\qquad n=m-r+2.\tag{*}\]
特别地，若 $n\geq\operatorname{prof}_s(F)+\operatorname{prof}_x^L(X_s)$，例如 $f$ 关于 $L$ \Term{locally-acyclic}{局部无环}，则
\[\operatorname{prof}_x(f^*F)\geq\operatorname{prof}_s(F)+\operatorname{prof}_x^L(X_s).\tag{**}\]
若 $L$ 只有一个元素 $\ell$，且 $n\geq\operatorname{prof}_s(F)+\operatorname{prof}_x^L(X_s)$，则上述不等式为等式。
\end{SGAresult}
\begin{proof}
取 $S$ 与 $x$ 在 $s$ 上方的\Term{geometric-point}{几何点} $\bar s$ 以及 $x$ 与 $\bar s$ 上方的\Term{geometric-point}{几何点} $\bar x$ 处的\Term{strict-localization}{严格局部化}，便归结到 $s,x$ 都为\Term{closed-point}{闭点}的情形。若 $n\geq\operatorname{prof}_s(F)+\operatorname{prof}_x^L(X_s)$，则在1.13中取 $p=\operatorname{prof}_s(F)$、$q=\operatorname{prof}_x^L(X_s)$ 即得(*)（$S-\{s\}$ 在 $S$ 中\Term{retrocompact}{逆紧}这一假设由 $X-X_s$ 在 $X$ 中\Term{retrocompact}{逆紧}及 $f$ 为\Term{surjectivity}{满射}得到；$f$ 因为 $(-1)$-\Term{acyclic}{无环}而为\Term{surjectivity}{满射}，除非1.15的结论为空）。若 $n<\operatorname{prof}_s(F)+\operatorname{prof}_x^L(X_s)$，也可由1.13得到(*)，例如取 $p=\operatorname{prof}_s(F)$、$q=n-p$。只须证明最后一个断言。置 $p=\operatorname{prof}_s(F)$、$q=\operatorname{prof}_x(X_s)$，由(1.13.2)有
\[\mathcal H_x^{p+q}(f^*(F))\simeq\mathcal H_x^q(f^*(\mathcal H_s^p(F))).\]
因 $\operatorname{prof}_s(F)=p$，$\mathcal H_s^p(F)$ 为非零的 $\ell$-\Term{torsion}{挠}\Term{sheaf}{层}，在 $s$ 上常值。因此 $G=f^*(\mathcal H_s^p(F))$ 为 $X_s$ 上常值的非零 $\ell$-\Term{torsion}{挠}\Term{sheaf}{层}，故包含一个\Term{isomorphism}{同构}于 $\mathbf Z/\ell\mathbf Z$ 的\Term{subsheaf}{子层}。由于 $\mathcal H_x^q(\mathbf Z/\ell\mathbf Z)\ne0$，确有 $\mathcal H_x^q(G)\ne0$。
\end{proof}
\begin{SGAresult}{推论}{1.16}
设 $f:X\to S$ 为\Term{characteristic}{特征}零的\Term{excellent}{优秀}\Term{scheme}{概形}之间的\Term{regular-morphism}{正则态射}（EGA IV 7.8.2），$\ell$ 为素数，$F$ 为 $S$ 上 $\ell$-\Term{torsion}{挠}\Term{sheaf}{层}的\Term{complex}{复形}。设 $x\in X$、$s=f(x)$，则
\[\operatorname{prof}_x(f^*F)=\operatorname{prof}_s(F)+2\dim(\mathcal O_{X,x}).\]
\end{SGAresult}
\begin{proof}
事实上 $f$ \Term{locally-acyclic}{局部无环}（SGA 4 XIX 4.1）。由1.15，
\[\operatorname{prof}_x(f^*F)=\operatorname{prof}_s(F)+\operatorname{prof}_x(X_s).\]
而由1.10，$\operatorname{prof}_x(X_s)=2\dim\mathcal O_{X_s,x}$，故得结论。
\end{proof}
\textbf{注记1.17。}由1.15可知，若将b)、c)换成如下条件，1.13仍成立：

b$'$) 对每点 $s\in f(Y)$，有 $\operatorname{prof}_s(F)\geq p$。

c$'$) 对每点 $x\in Y$，若 $s=f(x)$，则 $\operatorname{prof}_x^L(X_s)\geq q$。
\SourcePage{182}
对于\Term{sheaf-of-sets}{集合层}或\Term{sheaf-of-groups}{群层}，有如下类似1.13的定理。
\begin{SGAresult}{定理}{1.18}
设 $f:X\to S$ 为\Term{scheme}{概形}的\Term{morphism}{态射}，$Y$ 为 $X$ 的\Term{closed-subset}{闭子集}，$X-Y$ 在 $X$ 中\Term{retrocompact}{逆紧}，且对 $Y$ 的每点 $x$，\Term{local-ring}{局部环} $\mathcal O_{X,x}$ 为\Term{noetherian}{诺特}环。

$1^\circ)$ 设 $F$ 为 $S$ 上的\Term{sheaf-of-sets}{集合层}，$n=1$ 或2。假设 $f$ \Term{locally-n-acyclic}{局部 $(n-2)$-无环}，且对 $f(Y)$ 的每点 $s$ 有
\[\operatorname{prof}_{Y_s}(X_s)+\operatorname{prof}_s(F)\geq n.\]
则 $\operatorname{prof}_Y(f^*F)\geq n$。

$2^\circ)$ 设 $L$ 为素数\Term{set}{集合}，$F$ 为 $\operatorname{ind}$-$L$-群的\Term{sheaf}{层}。假设 $f$ 关于 $L$ \Term{locally-n-aspherical}{局部 $1$-非球面}（SGA 4 XV 1.11），且对 $f(Y)$ 的每点 $s$ 有
\[\operatorname{prof\,hop}_{Y_s}^L(X_s)+\operatorname{prof}_s(F)\geq3.\]
则 $\operatorname{prof}_Y(f^*F)\geq3$。
\end{SGAresult}
\begin{proof}
如1.14、1.15那样，归结到 $X,S$ 为\Term{strictly-local-scheme}{严格局部概形}，$f$ 为\Term{local-homomorphism}{局部同态}，$Y$ 为 $X$ 的\Term{closed-point}{闭点} $x$ 的情形。令 $s=f(x)$ 为 $S$ 的\Term{closed-point}{闭点}，有\Term{commutative-diagram}{交换图}：
\[\begin{tikzcd}
X-X_s\arrow[r,"i"]\arrow[d,"h"']&X-\{x\}\arrow[r,"j"]\arrow[dr,"g"]&X\arrow[d,"f"]\\
S-\{s\}\arrow[rr,"k"']&&S.
\end{tikzcd}\]
$1^\circ)$ a) $n=1$ 的情形。

若 $\operatorname{prof}_s(F)\geq1$，则\Term{morphism}{态射} $F\to k_*k^*F$ 为\Term{injectivity}{单射}，故 $f^*F\to f^*(k_*k^*F)$ 亦为\Term{injectivity}{单射}。又因 $f$ \Term{locally-n-acyclic}{局部 $(-1)$-无环}，\Term{morphism}{态射} $f^*(k_*k^*F)\to(j\cdot i)_*(f^*F|_{X-X_s})$ 为\Term{injectivity}{单射}。故复合\Term{morphism}{态射} $f^*F\to(j\cdot i)_*(f^*F|_{X-X_s})$ 为\Term{injectivity}{单射}，这说明 $\operatorname{prof}_{X_s}(f^*F)\geq1$，从而 $\operatorname{prof}_x(f^*F)\geq1$。

若 $\operatorname{prof}_x(X_s)\geq1$，考虑\Term{commutative-diagram}{交换图}
\[\begin{tikzcd}
H^0(X,f^*F)\arrow[r,"v"]\arrow[d,"\sim"]&H^0(X-\{x\},f^*F)\arrow[d]\\
H^0(X_s,f^*F)\arrow[r,"v'"']&H^0(X-\{x\},f^*F).
\end{tikzcd}\tag{*}\]
依假设，$v'$ 为\Term{injectivity}{单射}，故 $v$ 亦然。
\SourcePage{183}
b) $n=2$ 的情形。考虑\Term{commutative-diagram}{交换图}
\[\begin{tikzcd}[column sep=small]
H^0(S,F)\arrow[rr,"u"]\arrow[d,"m"',"\sim"]&&H^0(S-\{s\},F)\arrow[d,"n","\sim"']\\
H^0(X,f^*F)\arrow[r,"v"']&H^0(X-\{x\},f^*F)\arrow[r,"w"']&H^0(X-X_s,f^*F).
\end{tikzcd}\tag{**}\]
须说明 $v$ 为\Term{bijection}{双射}。$m$ 显然为\Term{bijection}{双射}，且 $f$ 为 $0$-\Term{acyclic}{无环}，故 $n$ 亦为\Term{bijection}{双射}。

若 $\operatorname{prof}_s(F)\geq2$，则 $u$ 为\Term{bijection}{双射}。如a)所见，仅由 $\operatorname{prof}_s(F)\geq1$ 已可推出 $\operatorname{prof}_{X_s}(f^*F)\geq1$；故 $v,w$ 为\Term{injectivity}{单射}，再由(**)可知 $v$ 为\Term{bijection}{双射}。

若 $\operatorname{prof}_x(X_s)\geq2$，则 $g$ 为 $0$-\Term{acyclic}{无环}（因为它\Term{locally-n-acyclic}{局部 $0$-无环}，且它的\Term{fiber}{纤维}为 $0$-\Term{acyclic}{无环}）。故 $v\cdot m$ 为\Term{bijection}{双射}，因而 $v$ 为\Term{bijection}{双射}。

若 $\operatorname{prof}_s(F)\geq1$ 且 $\operatorname{prof}_x(X_s)\geq1$，则已经知道 $v,w$ 为\Term{injectivity}{单射}。取 $X_s-\{x\}$ 的一个\Term{maximal-point}{极大点} $z$（由假设 $\operatorname{prof}_x(X_s)\geq1$，这样的点存在），令 $\overline Z$ 为 $X$ 在 $z$ 上方的一个\Term{geometric-point}{几何点}处的\Term{strict-localization}{严格局部化}，$\overline{f^*F}$ 为 $f^*F$ 在 $\overline Z$ 上的\Term{inverse-image}{逆像}。考虑\Term{commutative-diagram}{交换图}
\[\begin{tikzcd}[column sep=small]
H^0(S,F)\arrow[rr]\arrow[d,"m"',"\sim"]&&H^0(S-\{s\},F)\arrow[d,"n","\sim"']\\
H^0(X,f^*F)\arrow[r,"v"]\arrow[d,"m'"',"\sim"]&H^0(X-\{x\},f^*F)\arrow[r,"w"]\arrow[dl,"r"]&H^0(X-X_s,f^*F)\arrow[d,"n'","\sim"']\\
H^0(\overline Z,\overline{f^*F})\arrow[rr]&&H^0(\overline Z-\{\bar z\},\overline{f^*F}).
\end{tikzcd}\]
\Term{morphism}{态射} $m'\cdot m$ 显然为\Term{bijection}{双射}，而由 $f$ \Term{locally-n-acyclic}{局部 $0$-无环}，$n'\cdot n$ 为\Term{bijection}{双射}，故 $m',n'$ 亦为\Term{bijection}{双射}。由于 $w$ 为\Term{injectivity}{单射}，$r$ 亦为\Term{injectivity}{单射}，因而 $v$ 为\Term{bijection}{双射}。

$2^\circ)$ 由b)，已知 $\operatorname{prof}_x(f^*F)\geq2$。

若 $\operatorname{prof}_s(F)\geq3$，则 $R^1k_*(k^*F)=1$\footnote{编者注：后文先后用0或1表示\Term{trivial-bundle}{平凡}\Term{torsor}{挠子}；保留了这种双重记法，因为仔细考虑后可见它并不造成歧义。}。因 $f$ \Term{locally-n-aspherical}{局部 $1$-非球面}，
\[R^1(j\cdot i)_*(f^*F|_{X-X_s})=f^*(R^1k_*(k^*F))=1.\]
于是 $\operatorname{prof}_{X_s}(f^*F)\geq3$，从而 $\operatorname{prof}_x(f^*F)\geq3$。

若 $\operatorname{prof}_x(X_s)\geq3$，则 $g$ 为 $1$-\Term{aspherical}{非球面}（它\Term{locally-n-aspherical}{局部 $1$-非球面}，且它的\Term{fiber}{纤维}为 $1$-\Term{aspherical}{非球面}）。故 $H^1(X-\{x\},f^*F)=H^1(S,F)=1$，从而 $\operatorname{prof}_x(f^*F)\geq3$。

若 $\operatorname{prof}_s(F)\geq2$ 且 $\operatorname{prof}_x(X_s)\geq1$，则使用\Term{exact-sequence}{正合列}（SGA 4 XII 3.2）：
\[1\longrightarrow R^1j_*(i_*(f^*F|_{X-X_s}))\longrightarrow R^1(j\cdot i)_*(f^*F|_{X-X_s})\longrightarrow j_*(R^1i_*(f^*F|_{X-X_s})).\]
\SourcePage{184}
因 $f,g$ \Term{locally-n-aspherical}{局部 $1$-非球面}，有
\[\begin{aligned}
R^1(j\cdot i)_*(f^*F|_{X-X_s})&\simeq f^*(R^1k_*(k^*F)),\\
R^1i_*(f^*F|_{X-X_s})&\simeq g^*(R^1k_*(k^*F)).
\end{aligned}\]
故前述\Term{exact-sequence}{正合列}写成
\[1\longrightarrow R^1j_*(i_*(f^*F|_{X-X_s}))\longrightarrow f^*(R^1k_*(k^*F))\xrightarrow a j_*(j^*(f^*(R^1k_*(k^*F)))).\tag{***}\]
假设 $\operatorname{prof}_s(F)\geq2$ 表明 $F\to k_*k^*F$ 为\Term{bijection}{双射}；应用 $g^*$ 并利用 $g$ \Term{locally-n-acyclic}{局部 $0$-无环}，得到 $f^*F|_{X-\{x\}}=i_*(f^*F|_{X-X_s})$。假设 $\operatorname{prof}_x(X_s)\geq1$ 表明 $a$ 为\Term{injectivity}{单射}（注意 $f^*(R^1k_*(k^*F))$ 在 $X_s$ 外等于1，在 $X_s$ 上常值）。由(***)得 $R^1j_*(f^*F|_{X-\{x\}})=1$，故 $\operatorname{prof}_x(f^*F)\geq3$。

若 $\operatorname{prof}_s(F)\geq1$ 且 $\operatorname{prof}_x(f^*F)\geq2$，考虑由\Term{exact-sequence}{正合列}
\[1\longrightarrow F\longrightarrow k_*k^*F\longrightarrow G\longrightarrow1\]
定义的\Term{homogeneous-space}{齐性空间}\Term{sheaf}{层} $G$。对这个\Term{exact-sequence}{正合列}应用\Term{exact-functor}{正合函子} $g^*$，并利用（SGA 4 XII 3.1），得到如下各行\Term{exact}{正合}的\Term{commutative-diagram}{交换图}：
\[\begin{tikzcd}[column sep=small]
f^*(k_*k^*F)\arrow[r]\arrow[d]&f^*G\arrow[r]\arrow[d,"b"]&1&{}\\
j_*(g^*(k_*k^*F))\arrow[r,"u"]&j_*(g^*G)\arrow[r]&R^1j_*(g^*F)\arrow[r]&R^1j_*(g^*(k_*k^*F)).
\end{tikzcd}\]
因 $\operatorname{prof}_x(X_s)\geq2$，$b$ 为\Term{bijection}{双射}，故 $u$ 为\Term{surjectivity}{满射}，从而有如下核仅含中性元的映射：
\[1\longrightarrow R^1j_*(g^*F)\longrightarrow R^1j_*(g^*(k_*k^*F))=R.\]
由于 $g^*(k_*k^*F)\simeq i_*(f^*F|_{X-X_s})$（因 $g$ \Term{locally-n-acyclic}{局部 $0$-无环}），$R$ 等同于\Term{exact-sequence}{正合列}(***)的首项。而前一种情形已说明，只要 $\operatorname{prof}_x(X_s)\geq1$，就有 $R=1$。故 $\operatorname{prof}_x(f^*F)\geq3$，1.18的证明完成。
\end{proof}
下面的推论推广（SGA 4 XVI 3.2和3.3）。
\begin{SGAresult}{推论}{1.19}
设 $f:X\to S$ 为\Term{locally-noetherian}{局部诺特}\Term{scheme}{概形}之间的\Term{flat-morphism}{平坦态射}，其\Term{fiber}{纤维}\Term{separable}{可分}；$Y$ 为 $X$ 的\Term{closed-subset}{闭子集}。假设对每点 $s\in f(Y)$，\Term{fiber}{纤维} $Y_s$ 在 $X_s$ 中\Term{rare}{稀疏}\footnote{编者注：“\Term{rare}{稀疏}”即“内部为空”，参看 Bourbaki TG IX.52。}，且下列条件之一成立：

a) $f(Y)$ 的\Term{closure}{闭包}在 $S$ 中\Term{rare}{稀疏}。

b) $X_s$ 在 $Y_s$ 的各点\Term{geometrically-unibranch}{几何单分支}。

则 $\operatorname{prof}_Y(X)\geq2$。
\end{SGAresult}
\SourcePage{185}
\begin{proof}
由关于 $f$ 的假设，$f$ \Term{locally-n-acyclic}{局部 $0$-无环}（SGA 4 XV 4.1）。应用1.13。由1.6 b)，假设 $Y_s$ 在 $X_s$ 中\Term{rare}{稀疏}（分别 $f(Y)$ 在 $S$ 中\Term{rare}{稀疏}）\Term{equivalence}{等价}于 $\operatorname{prof}_{Y_s}(X_s)\geq1$（分别 $\operatorname{prof}_{\overline{f(Y)}}(S)\geq1$）。$X_s$ 在 $Y_s$ 的每点\Term{geometrically-unibranch}{几何单分支}\Term{equivalence}{等价}于 $X_s$ 在 $Y_s$ 的\Term{geometric-point}{几何点}处的\Term{strict-localization}{严格局部化}为\Term{irreducible}{不可约}。因 $Y_s$ 在 $X_s$ 中\Term{rare}{稀疏}，由1.8这显然蕴涵 $\operatorname{prof}_{Y_s}(X_s)\geq2$。无论哪种情形，1.13都给出 $\operatorname{prof}_Y(X)\geq2$。
\end{proof}
\begin{SGAresult}{推论}{1.20}
设 $f:X\to S$ 为\Term{locally-noetherian}{局部诺特}\Term{scheme}{概形}之间的\Term{regular-morphism}{正则态射}（EGA IV 6.8.1），$Y$ 为 $X$ 的\Term{closed-subset}{闭子集}。假设对每点 $s\in f(Y)$，下列条件之一成立：

a) $\operatorname{codim}(Y_s,X_s)\geq2$。

b) $\operatorname{codim}(Y_s,X_s)\geq1$ 且 $\operatorname{prof}_s(S)\geq1$。

c) $\operatorname{prof\,hop}_s(S)\geq3$。

则 $\operatorname{prof\,hop}_Y(X)\geq3$。
\end{SGAresult}
\begin{proof}
这由1.18得到，因为假设a)蕴涵 $\operatorname{prof\,hop}_{Y_s}(X_s)\geq3$（参看1.11），而 $\operatorname{codim}(Y_s,X_s)\geq1$ 显然蕴涵 $\operatorname{prof}_Y(X)\geq2$。
\end{proof}
\SGAsection{2}{技术性引理}
\textbf{2.1。}设 $S$ 为\Term{locally-noetherian}{局部诺特}\Term{scheme}{概形}，$f:X\to S$ 为\Term{locally-finite-type}{局部有限型}\Term{morphism}{态射}，$t\in S$。若 $x\in X$ 满足 $s=f(x)\in\operatorname{Spec}\mathcal O_{S,t}$，置
\[\delta_t(x)=\operatorname{deg\,tr}k(x)/k(s)+\dim(\overline{\{s\}}),\]
其中 $\overline{\{s\}}$ 为 $s$ 在 $\operatorname{Spec}\mathcal O_{S,t}$ 中的\Term{closure}{闭包}，$k(x),k(s)$ 分别为 $x,s$ 的\Term{residue-field}{剩余域}。若 $S$ 为\Term{closed-point}{闭点}为 $t$ 的\Term{local-ring}{局部环}，也以 $\delta(x)$ 代替 $\delta_t(x)$（参看 SGA 4 XIV 2.2）。
\begin{SGAresult}{引理}{2.1.1}
设有\Term{cartesian-square}{笛卡尔方块}
\[\begin{tikzcd}X'\arrow[r,"h"]\arrow[d,"f'"']&X\arrow[d,"f"]\\S'\arrow[r,"g"']&S,\end{tikzcd}\]
其中 $S,S'$ 为\Term{noetherian}{诺特}\Term{local-ring}{局部环}，\Term{closed-point}{闭点}分别为 $t,t'$；$g$ 为\Term{faithfully-flat}{忠实平坦}\Term{morphism}{态射}，满足 $g^{-1}(t)=t'$；$f$ 为\Term{locally-finite-type}{局部有限型}\Term{morphism}{态射}。设 $x'\in X'$、$x=h(x')$、$s=f(x)$、$s'=f'(x')$，则
\[\delta(x')\leq\delta(x).\]
\SourcePage{186}
此外，上述不等式是不等式，当且仅当
\[\operatorname{deg\,tr}k(x)/k(s)=\operatorname{deg\,tr}k(x')/k(s')\quad\text{且}\quad\dim(\overline{\{s\}})=\dim(\overline{\{s'\}}).\]
特别地，给定 $x\in X$，可找到 $x'$ 使 $\delta(x)=\delta(x')$。
\end{SGAresult}
\begin{proof}
事实上有（EGA IV 6.11）$\dim(\overline{\{s\}})=\dim g^{-1}(\overline{\{s\}})$。故对 $g^{-1}(s)$ 的每点 $s'$，有 $\dim(\overline{\{s'\}})\leq\dim(\overline{\{s\}})$，而给定 $s$ 后可找到 $s'\in g^{-1}(s)$ 使等号成立。记 $Z$ 为 $x$ 在 $X$ 于 $s$ 的\Term{fiber}{纤维} $X_s$ 中的\Term{schematic-closure}{概形论闭包}，令 $Z'=Z\times_{\operatorname{Spec}k(s)}\operatorname{Spec}k(s')$。则 $Z'$ 为\Term{equidimensional}{等维}的，\Term{dimension}{维数}为 $\operatorname{deg\,tr}k(x)/k(s)$。故对每点 $x'\in Z'_x$ 有
\[\operatorname{deg\,tr}k(x')/k(s')\leq\operatorname{deg\,tr}k(x)/k(s),\]
当 $x'$ 为 $Z'_x$ 的\Term{maximal-point}{极大点}时等号成立。所述结论随即得到。
\end{proof}
\textbf{2.2。}设 $f:X\to S$ 为\Term{locally-finite-type}{局部有限型}\Term{morphism}{态射}，$T$ 为 $S$ 的\Term{closed-subset}{闭子集}。设 $x\in X$、$s=f(x)$，置
\[\delta_T(x)=\operatorname{deg\,tr}k(x)/k(s)+\operatorname{codim}(\overline{\{s\}}\cap T,\overline{\{s\}})=\inf_{t\in T\cap\overline{\{s\}}}\delta_t(x).\]
\begin{SGAresult}{引理}{2.2.1}
设有\Term{cartesian-square}{笛卡尔方块}
\[\begin{tikzcd}X'\arrow[r,"h"]\arrow[d,"f'"']&X\arrow[d,"f"]\\S'\arrow[r,"g"']&S,\end{tikzcd}\]
其中\Term{scheme}{概形} $S,S'$ 为\Term{locally-noetherian}{局部诺特}且\Term{catenary}{链状}，\Term{morphism}{态射} $f$ 为\Term{locally-finite-type}{局部有限型}，$g$ 为\Term{faithfully-flat}{忠实平坦}。设 $T$ 为 $S$ 的\Term{closed-subset}{闭子集}，$T'$ 为 $S'$ 的\Term{closed-subset}{闭子集}，$g(T')\subset T$；$x'\in X'$、$x=h(x')$，并以
\[h_{x'}:\operatorname{Spec}\mathcal O_{X',x'}\longrightarrow\operatorname{Spec}\mathcal O_{X,x}\]
记 $h$ 诱导的\Term{morphism}{态射}。则
\[\delta_T(x)-\delta_{T'}(x')\leq\dim h_{x'}^{-1}(x).\]
\end{SGAresult}
\begin{proof}
设 $s'=f'(x')$、$s=f(x)$。依定义，
\[\begin{aligned}
\delta_T(x)-\delta_{T'}(x')={}&\operatorname{deg\,tr}k(x)/k(s)-\operatorname{deg\,tr}k(x')/k(s')\\
&+\operatorname{codim}(\overline{\{s\}}\cap T,\overline{\{s\}})-\operatorname{codim}(\overline{\{s'\}}\cap T',\overline{\{s'\}}).
\end{aligned}\]
因 $g$ \Term{faithfully-flat}{忠实平坦}，由（EGA IV 6.1.4），
\[\begin{aligned}
\operatorname{codim}(\overline{\{s\}}\cap T,\overline{\{s\}})&=\operatorname{codim}(g^{-1}(\overline{\{s\}})\cap g^{-1}(T),g^{-1}(\overline{\{s\}}))\\
&\leq\operatorname{codim}(g^{-1}(\overline{\{s\}})\cap T',g^{-1}(\overline{\{s\}})).
\end{aligned}\tag{*}\]
\SourcePage{187}
因 $S'$ \Term{catenary}{链状}，由（EGA $0_{\mathrm{IV}}$ 14.3.2 b)），
\[\begin{aligned}
\operatorname{codim}(\overline{\{s'\}}\cap T',g^{-1}(\overline{\{s\}}))
&=\operatorname{codim}(\overline{\{s'\}}\cap T',\overline{\{s'\}})+\operatorname{codim}(\overline{\{s'\}},g^{-1}(\overline{\{s\}}))\\
&=\operatorname{codim}(\overline{\{s'\}}\cap T',g^{-1}(\overline{\{s\}})\cap T')\\
&\quad+\operatorname{codim}(g^{-1}(\overline{\{s\}})\cap T',g^{-1}(\overline{\{s\}})).
\end{aligned}\]
由此及(*)推出
\[\begin{aligned}
\delta_T(x)-\delta_{T'}(x')\leq{}&\operatorname{deg\,tr}k(x)/k(s)-\operatorname{deg\,tr}k(x')/k(s')\\
&+\operatorname{codim}(\overline{\{s'\}},g^{-1}(\overline{\{s\}})).
\end{aligned}\]
计算 $\operatorname{codim}(\overline{\{s'\}},g^{-1}(\overline{\{s\}}))=\dim\mathcal O_{S'_s,s'}$（其中 $S'_s$ 为 $S'$ 在 $s$ 的\Term{fiber}{纤维}）。设 $Z$ 为 $x$ 在 $X_s$ 中的闭像，$Z'\subset X'_s$ 为由\Term{cartesian-square}{笛卡尔方块}
\[\begin{tikzcd}Z'\arrow[r]\arrow[d]&Z\arrow[d]\\S'_s\arrow[r]&\operatorname{Spec}k(s)\end{tikzcd}\]
定义的\Term{scheme}{概形}。\Term{morphism}{态射} $Z\to\operatorname{Spec}k(s)$ 为\Term{flat}{平坦}、\Term{locally-finite-type}{局部有限型}，且 $\dim Z=\operatorname{deg\,tr}k(x)/k(s)$。由（EGA IV 6.1.2）有
\[\dim\mathcal O_{Z',x'}=\dim\mathcal O_{S'_s,s'}+\operatorname{deg\,tr}k(x)/k(s)-\operatorname{deg\,tr}k(x')/k(s').\]
结合 $Z'_{s'}\simeq Z\otimes_{k(s)}k(s')$，得到
\[\delta_T(x)-\delta_{T'}(x')\leq\dim\mathcal O_{Z',x'}.\]
而 $\operatorname{Spec}(\mathcal O_{Z',x'})$ 等同于\Term{morphism}{态射}
\[(\operatorname{Spec}\mathcal O_{X',x'})_s\longrightarrow(\operatorname{Spec}\mathcal O_{X,x})_s\]
在 $x$ 的\Term{fiber}{纤维}，从而也等同于 $h_{x'}$ 在 $x$ 的\Term{fiber}{纤维}，这就证明了定理。
\end{proof}
\textbf{2.3。}第4节各定理的证明以\Term{duality-theory}{对偶理论}为基础，使用下列引理。设整数 $m$ 是素数 $\ell$ 的幂；若 $X$ 为\Term{scheme}{概形}，则所考虑的 $X$ 上的层一律为 $\mathbf Z/m\mathbf Z$-\Term{module-sheaf}{模层}，于是有 $X$ 上的\Term{dualizing-complex}{对偶化复形}的概念（SGA 5 I 1.7）。假设 $X$ 上存在这样的\Term{complex}{复形} $K$；则对 $X$ 的每点 $x$ 上方的\Term{geometric-point}{几何点} $\bar x$，由 $K$ 得到（参看 SGA 5 I 4.5）$\operatorname{Spec}k(\bar x)$ 上的\Term{dualizing-complex}{对偶化复形} $K_{\bar x}$，且 $K_{\bar x}\simeq\mathbf Z/m\mathbf Z[n]$（方括号表示\Term{translation-functor}{平移函子}），其中整数 $n$ 只依赖于 $x$。置
\[\delta^K(x)=n.\]
故若 $K$ 在点 $x$ 处\Term{normalized}{归一化}（SGA 5 I 4.5），则 $n=0$。
\begin{SGAresult}{引理}{2.3.1}
设 $X$ 为配备\Term{dualizing-complex}{对偶化复形} $K$ 的\Term{locally-noetherian}{局部诺特}\Term{scheme}{概形}。若 $x,x'$ 为 $X$ 的两点，$x$ 是 $x'$ 的\Term{specialization}{特化}，且 $\operatorname{codim}(\overline{\{x\}},\overline{\{x'\}})=1$，则
\[\delta^K(x)=\delta^K(x')-2.\]
\end{SGAresult}
\SourcePage{188}
\begin{proof}
先可归结到 $X$ 为\Term{strictly-local-scheme}{严格局部概形}的情形。事实上，令 $\overline X$ 为 $X$ 在 $x$ 上方的\Term{geometric-point}{几何点} $\bar x$ 处的\Term{strict-localization}{严格局部化}，$i:\overline X\to X$ 为\Term{canonical}{典范}\Term{morphism}{态射}，$\bar x'$ 为 $\overline X$ 中 $x'$ 上方的\Term{geometric-point}{几何点}。则 $i^*K$ 为 $\overline X$ 上的\Term{dualizing-complex}{对偶化复形}，且有（SGA 5 I 4.5）
\[(i^*K)_{\bar x}\simeq K_{\bar x},\qquad(i^*K)_{\bar x'}\simeq K_{\bar x'},\]
从而完成对\Term{strictly-local}{严格局部}情形的归结。

若 $j:\overline{\{x'\}}\to X$ 为 $X$ 的底空间为 $\overline{\{x'\}}$ 的\Term{reduced}{既约}\Term{closed-subscheme}{闭子概形}的\Term{immersion}{浸入}，则 $R^!j(K)$ 是 $\overline{\{x'\}}$ 上的\Term{dualizing-complex}{对偶化复形}。由（SGA 5 I 4.5）立即看出，只须对 $\overline{\{x'\}}$ 证明引理。故归结到 $X$ 为一维\Term{integral}{整}\Term{strictly-local-scheme}{严格局部概形}的情形。

令 $X'$ 为 $X$ 的\Term{normalization}{正规化}，$f:X'\to X$ 为\Term{canonical}{典范}\Term{morphism}{态射}；$f$ 是\Term{integral-morphism}{整}、\Term{surjectivity}{满射}、\Term{radicial}{根性}\Term{morphism}{态射}，故 $f^*K$ 是 $X'$ 上的\Term{dualizing-complex}{对偶化复形}，且只须对 $X'$ 以及 $x,x'$ 上方的点证明引理。于是归结到 $X$ 为一维\Term{local-scheme}{局部}、\Term{integral}{整}、\Term{regular}{正则}\Term{scheme}{概形}的情形。但已知（参看 SGA 5 I 4.6.2和5.1）$\mu_m[2]$ 与 $\mathbf Z/m\mathbf Z$ 此时为分别在 $x,x'$ 处\Term{normalized}{归一化}的\Term{dualizing-complex}{对偶化复形}，引理随即得到。
\end{proof}
\begin{SGAresult}{引理}{2.3.2}
设 $S$ 为\Term{noetherian}{诺特}\Term{local-scheme}{局部概形}，$f:X\to S$ 为\Term{finite-type}{有限型}\Term{morphism}{态射}。若 $K$ 是 $S$ 上在其\Term{closed-point}{闭点} $t$ 处\Term{normalized}{归一化}的\Term{dualizing-complex}{对偶化复形}，且 $R^!f(K)=K'$ 为 $X$ 上的\Term{dualizing-complex}{对偶化复形}（参看 SGA 5 I 3.4.3），则对 $X$ 的每点 $x$，
\[\delta^{K'}(x)=2\delta(x).\]
\end{SGAresult}
\begin{proof}
令 $s=f(x)$，$x'$ 为\Term{fiber}{纤维} $X_s$ 的一个\Term{closed-point}{闭点}，则 $\delta^{K'}(x')=\delta^K(s)$，并由2.3.1，
\[\delta^K(s)=2\operatorname{codim}(\overline{\{t\}},\overline{\{s\}})=2\dim(\overline{\{s\}}).\]
可选 $x'$ 为 $x$ 的\Term{specialization}{特化}，故由2.3.1，
\[\delta^{K'}(x)=\delta^{K'}(x')+2\operatorname{codim}(\overline{\{x\}},\overline{\{x'\}})=\delta^{K'}(x')+2\operatorname{deg\,tr}k(x)/k(s).\]
引理随即得到。
\end{proof}
下面的引理只在第4节 Lefschetz 定理的逆命题中使用。
\begin{SGAresult}{引理}{2.3.3}
设有\Term{cartesian-square}{笛卡尔方块}
\[\begin{tikzcd}X'\arrow[r,"h"]\arrow[d,"f'"']&X\arrow[d,"f"]\\S'\arrow[r,"g"']&S,\end{tikzcd}\]
其中 $S$ 为\Term{characteristic}{特征}零的\Term{excellent}{优秀}\Term{strictly-local-scheme}{严格局部概形}，$S'$ 为 $S$ 的\Term{completion}{完备化}，$f$ 为\Term{finite-type}{有限型}\Term{morphism}{态射}。设 $\ell$ 为素数，$x\in X$，$Z$ 为 $X'_x$ 在 $X'$ 中的\Term{schematic-closure}{概形论闭包}，
\SourcePage{189}
$i:X'_x\to Z$、$j:Z\to X$ 为\Term{canonical}{典范}\Term{morphism}{态射}。则若 $k:X'\to R$ 为 $X'$ 在\Term{characteristic}{特征}零的\Term{excellent}{优秀}\Term{regular}{正则}\Term{scheme}{概形} $R$ 中的\Term{closed-immersion}{闭浸入}，\Term{complex}{复形}
\[K'=i^*(R^!(k\cdot j)(\mathbf Z/\ell\mathbf Z))\]
为 $X'_x$ 上的常值\Term{dualizing-complex}{对偶化复形}（即只有一个非零\Term{cohomology-sheaf}{上同调层}，且其\Term{isomorphism}{同构}于 $\mathbf Z/\ell\mathbf Z$）。
\end{SGAresult}
\begin{proof}
由（SGA 5 I 3.4.3），只须证明 $K'$ 为常值。由于 $Z$ \Term{excellent}{优秀}，$Z$ 中\Term{local-ring}{局部环}为\Term{regular}{正则}环的点的\Term{set}{集合}是一个\Term{open-subset}{开集} $U$（EGA IV 7.8.3 (iv)），且 $U$ 显然包含\Term{regular}{正则}的 $X'_x$。令 $u:U\to R$ 为 $U$ 在 $R$ 中的\Term{canonical}{典范}\Term{immersion}{浸入}，则由\Term{purity-theorem}{纯性定理}（SGA 4 XIX 3.2和3.4）及\Term{isomorphism}{同构}
\[(\mu_\ell)_S\simeq(\mathbf Z/\ell\mathbf Z)_S\]
（$S$ 为\Term{strictly-local-scheme}{严格局部概形}），得到
\[R^!u(\mathbf Z/\ell\mathbf Z)\simeq\mathbf Z/\ell\mathbf Z[2c],\]
其中 $c$ 为 $U$ 上的\Term{locally-constant}{局部常值}函数；它必在 $X'_x$ 的一个邻域内为常值，因为由（EGA IV 18.9.1），$g$ 的\Term{fiber}{纤维}是\Term{geometrically-integral}{几何整}的，故 $X'_x$ 为\Term{integral}{整}的。引理随即得到。
\end{proof}
